\chapter{Real-time vs Batch Inference (Research and Deployment)}

\ChapterMeta{
This chapter compares batch and real-time operating modes and outlines the canonical TensorRT pipeline with parity and benchmark safeguards.
}{
It guides you to high-throughput deployment paths while keeping artifact contracts and drift detection in place.
}{
Real-time TensorRT paths are typically Linux with an NVIDIA GPU; parity requires consistent preprocessing and output formats, and expensive extras like TTT and refinement should be explicitly budgeted.
}
\ChapterAudience{Read this chapter when you need to choose between offline throughput and real-time deployment constraints.}


\section{Two operating modes}
\begin{description}
  \item[Batch/offline] Run over a dataset split, export \path{predictions.json}, then evaluate and plot.
  \item[Real-time/streaming] Optimize end-to-end latency and stability; avoid expensive per-frame extras.
\end{description}

\begin{figure}[h]
\centering
\resizebox{\linewidth}{!}{%
\begin{tikzpicture}[
  node distance=10mm,
  font=\small,
  box/.style={draw, rounded corners, align=left, inner sep=6pt},
  arrow/.style={-{Latex[length=2mm]}, thick},
  opt/.style={draw, rounded corners, align=left, inner sep=6pt, dashed}
]

\node[box] (in) {\textbf{Input}\\frame / batch};
\node[box, right=16mm of in] (pre) {\textbf{Preprocess}\\resize/letterbox\\normalize};
\node[opt, right=16mm of pre] (ttt) {\textbf{TTT} (optional)\\bounded-cost};
\node[box, right=16mm of ttt] (eng) {\textbf{Engine}\\torch / onnxrt / trt};
\node[box, right=16mm of eng] (post) {\textbf{Postprocess}\\decode / NMS\\format convert};
\node[opt, right=16mm of post] (ref) {\textbf{Refine / gates} (optional)\\constraints / symmetry};
\node[box, right=16mm of ref] (out) {\textbf{Emit}\\\path{predictions.json}\\+ reports};

\draw[arrow] (in) -- (pre);
\draw[arrow] (pre) -- (ttt);
\draw[arrow] (ttt) -- (eng);
\draw[arrow] (eng) -- (post);
\draw[arrow] (post) -- (ref);
\draw[arrow] (ref) -- (out);

\end{tikzpicture}
}
\caption{Real-time vs batch inference share the same artifact boundary (\cmd{predictions.json}); optional extras must be explicitly budgeted.}
\end{figure}

\section{Non-goals (operational boundaries)}
\begin{itemize}
  \item TTT is \textbf{not} a real-time default. Keep it disabled unless explicitly budgeted and enabled.
  \item If TTT is enabled in streaming mode, use bounded-cost knobs first (e.g., \cmd{--ttt-batch-size 1 --ttt-max-batches 1}).
  \item Hessian refinement is not part of the TensorRT inference graph; treat it as an optional postprocess step.
  \item Real-time claims without hardware-pinned latency/FPS evidence are out of scope for this chapter.
\end{itemize}

\section{Backend choices}
\begin{itemize}
  \item \textbf{Precomputed predictions}: best for integration; works on macOS; use \path{tools/run_scenarios.py} with \cmd{precomputed} adapter.
  \item \textbf{Torch}: flexible for research (TTT, ablations), but slower and less deployable.
  \item \textbf{ONNX Runtime}: CPU-friendly parity reference; useful in CI.
  \item \textbf{TensorRT}: preferred for real-time throughput on Linux+NVIDIA.
\end{itemize}

References:
\begin{itemize}
  \item \path{docs/real_model_interface.md}
  \item \path{docs/tensorrt_pipeline.md}
\end{itemize}

\section{TensorRT pipeline (canonical)}
The typical export route is PyTorch $\rightarrow$ ONNX $\rightarrow$ TensorRT (engine build), with parity checks.

Canonical pattern (artifacts-first):
\begin{itemize}
  \item Export ONNX (+ metadata JSON): \\
    \cmd{python3 tools/export_trt.py} \\
    \cmd{--onnx ...} \cmd{--opset 18} \cmd{--dynamic-hw ...}
  \item Build engine (+ metadata JSON): \\
    \cmd{python3 tools/build_trt_engine.py} \\
    \cmd{--onnx ...} \cmd{--engine ...} \cmd{--precision fp16}
  \item Export \path{predictions.json} from TensorRT and run parity vs ONNXRuntime predictions.
\end{itemize}

The main rule for parity is to keep preprocessing and output formats identical (letterbox/RGB, box format, score scaling), and keep NMS disabled in exported graphs when comparing raw outputs.

\noindent\textbf{Input shape note:} Unlike PyTorch eager mode, ONNXRuntime and TensorRT are typically \textbf{shape-pinned} unless you explicitly export with dynamic axes / optimization profiles. If an ONNX or engine expects (for example) $64\times64$, feeding $640\times640$ will fail. The TensorRT exporter defaults to inferring \cmd{imgsz} from the engine input when possible (and accepts an explicit \cmd{--imgsz} override for dynamic engines).

\section{Stage 7: Performance + deployment safeguards (implemented tooling)}
Stage 7 is about keeping the real-time path fast and predictable.
In this repo, the main safeguards are implemented as \textbf{tooling and contracts} rather than as model-internal graph logic.

\subsection{Keep symmetry/commonsense logic out of TensorRT graphs}
The TensorRT export/build pipeline is designed to focus on the model forward pass.
Symmetry, template scoring, commonsense constraints, and gates are intended to run in \textbf{postprocess} where possible.

\begin{figure}[h]
\centering
\resizebox{0.95\linewidth}{!}{%
\begin{tikzpicture}[
  node distance=10mm,
  font=\small,
  box/.style={draw, rounded corners, align=left, inner sep=6pt},
  group/.style={draw, rounded corners, inner sep=10pt},
  arrow/.style={-{Latex[length=2mm]}, thick}
]

\node[box] (pre) {\textbf{Preprocess}\\resize/letterbox\\normalize};
\node[box, right=18mm of pre] (fwd) {\textbf{Model forward}\\(TensorRT engine)};
\node[box, right=18mm of fwd] (post) {\textbf{Postprocess}\\decode / NMS\\format};
\node[box, below=10mm of post] (logic) {\textbf{Keep outside engine}\\constraints / symmetry\\TTT / refinement};

\draw[arrow] (pre) -- (fwd);
\draw[arrow] (fwd) -- (post);
\draw[arrow] (post) -- (logic);

\node[group, fit=(fwd), label={[align=left]above:TRT graph}] {};
\node[group, fit=(pre) (post) (logic), label={[align=left]above left:Python/C++ glue}] {};

\end{tikzpicture}
}
\caption{Boundary rule: keep non-forward logic out of the TensorRT graph to simplify deployment and make drift diagnosable.}
\end{figure}

In particular:
\begin{itemize}
  \item Commonsense constraints live in \path{yolozu/constraints.py} and are toggled via \path{configs/runtime/constraints.yaml}.
  \item Parity comparisons are performed on raw outputs (no-NMS exports) so numerical mismatches are detectable.
\end{itemize}

\subsection{Benchmark and regression surfaces}
There are two complementary benchmark surfaces:
\begin{itemize}
  \item A backend-agnostic latency/FPS harness: \path{tools/benchmark_latency.py}
  \item A TensorRT engine latency measurement tool (Linux+NVIDIA + TRT python): \path{tools/measure_trt_latency.py}
\end{itemize}

The harness supports optional targets and can fail the run if targets are not met:
\begin{lstlisting}[language=bash]
python tools/benchmark_latency.py --config configs/benchmark_latency_example.json \
  --target-fps 30
\end{lstlisting}

For a full end-to-end TensorRT run on a Linux+NVIDIA box, use:\\
\path{tools/run_trt_pipeline.py} (export/predictions/parity/eval/latency/benchmark; supports \cmd{--dry-run} for CI-style artifact checks).

\subsection{Gate checks against a baseline}
The repo includes a lightweight gate-check script that compares a current scenario suite report with a saved baseline:
\begin{itemize}
  \item \path{tools/check_gates.py} (writes \path{reports/gate_check.json})
  \item baseline source: \path{reports/baseline.json} (generated by \path{tools/run_baseline.py})
\end{itemize}

The current Stage 7 check is a simple \cmd{fps >= 30} threshold in the gate output.
Note that real FPS validation is inherently \textbf{hardware-dependent}; CI runs can validate the harness and contracts,
while deployment-relevant benchmarks should be run on reference Linux+NVIDIA hardware.

\section{Real-time constraints and optional features}
\subsection{TTT in real-time}
TTT can improve robustness under domain shift but adds latency.
If you use it in streaming mode, keep it bounded:
\begin{itemize}
  \item prefer \cmd{--ttt-preset safe}
  \item start with \cmd{--ttt-batch-size 1 --ttt-max-batches 1}
  \item prefer \cmd{--ttt-reset stream} when you accept cross-image state
\end{itemize}

For controlled comparisons and ablations, use \cmd{--ttt-reset sample} even though it is slower.

\subsection{Per-detection refinement}
Hessian refinement (depth/rotation/offsets) is usually a \textbf{post-inference research tool}.
It is not typically enabled in production unless the added cost is explicitly budgeted.

\section{Batch throughput tips}
\begin{itemize}
  \item Increase batch size first (until memory limits).
  \item Keep dataset subsets deterministic for comparisons.
  \item Record engine metadata (GPU/driver/TensorRT versions) when benchmarking.
\end{itemize}

\section{Mac development vs GPU execution}
If you develop on macOS, keep GPU/TensorRT runs on a Linux+NVIDIA box.
The repo provides a Runpod-oriented path; see \path{deploy/runpod/README.md}.

\subsection{Preflight split for GPU sweep tasks}
Before reserving GPU runtime, generate a local preflight split report:
\begin{lstlisting}[language=bash]
python3 tools/gpu_validation_preflight.py --output reports/gpu_validation_preflight.json
\end{lstlisting}

The report separates:
\begin{itemize}
  \item \path{summary.local_executable_steps}: tasks you can complete locally
  \item \path{summary.gpu_execution_required_steps}: tasks that require Linux+NVIDIA runtime
  \item \path{summary.needs_gpu_runtime_steps}: tasks ready to run once GPU runtime is available
\end{itemize}

Operational runbook: \path{docs/runpod_gpu_validation_split.md}.

When operating via GitHub Actions machine runners instead of RunPod sessions, use
\path{.github/workflows/gpu_zisn_pipeline.yml} (\cmd{workflow_dispatch}) and select
stage \cmd{zisn1}, \cmd{zisn2}, \cmd{zisn3}, or \cmd{all}. The workflow publishes deterministic
DoD summaries to \path{ci-logs/gpu-zisn}.

\section{Inference profiler integration}
The \texttt{yolozu.inference.profiler} module wraps \texttt{torch.profiler} to produce kernel-level latency traces that complement wall-clock benchmarks.

\begin{itemize}
  \item \texttt{profiler\_available()}\\
  checks for \texttt{torch.profiler.profile}.
  \item \texttt{profile\_inference(adapter, records, warmup, active, output\_dir)}\\
  profiles \texttt{adapter.predict()} calls and writes TensorBoard/Chrome-trace output.
  \item \texttt{profile\_callable(fn, *args, iterations, warmup)}\\
  profiles any callable and returns a text summary table.
\end{itemize}

\begin{lstlisting}[language=Python]
from yolozu.inference.profiler import profile_inference, profiler_available

if profiler_available():
    summary = profile_inference(adapter, records, output_dir="runs/profile")
    print(summary)
\end{lstlisting}

\section{torch.compile and ONNX export helpers}
The \texttt{yolozu.inference.torch\_export} module provides two orthogonal utilities:

\begin{itemize}
  \item \texttt{compile\_for\_inference(model, backend, mode)} --- wraps \texttt{torch.compile} with configurable backend (\texttt{inductor}, \texttt{aot\_eager}, etc.) and mode (\texttt{reduce-overhead}, \texttt{max-autotune}). Falls back gracefully when compilation is unavailable.
  \item \texttt{export\_model\_onnx(model, sample\_input, output\_path, opset\_version, ...)} --- exports a \texttt{torch.nn.Module} to ONNX format with configurable opset, dynamic axes, and named I/O.
  \item \texttt{torch\_compile\_available()}, \texttt{torch\_export\_available()} --- feature-detection guards.
\end{itemize}

The \texttt{RTDETRPoseAdapter} integrates \texttt{compile\_for\_inference} via its \texttt{compile\_model}, \texttt{compile\_backend}, and \texttt{compile\_mode} constructor kwargs; when enabled, the model is automatically compiled at first \texttt{predict()} call.

\begin{lstlisting}[language=Python]
from yolozu.inference.torch_export import (
    compile_for_inference, export_model_onnx,
)

fast = compile_for_inference(model, mode="reduce-overhead")
export_model_onnx(model, sample_input, "model.onnx", opset_version=17)
\end{lstlisting}

\section{What to report in papers/notes}
For real-time or deployment-relevant results, always report:
\begin{itemize}
  \item hardware (GPU, driver, CUDA/TensorRT) and engine build settings
  \item end-to-end latency / FPS and measurement method
  \item whether TTT/refinement/gates were enabled and their bounds
\end{itemize}
